\documentclass[11pt,a4paper]{article}
\usepackage[utf8]{inputenc}
\usepackage[vietnamese]{babel}
\usepackage{amsmath,amssymb}
\usepackage{geometry}
\geometry{a4paper, total={180mm,260mm}, left=15mm, right=15mm, top=18mm, bottom=18mm}

\usepackage{tikz}
\usetikzlibrary{shapes.geometric, arrows.meta, positioning, calc}
\tikzset{
  arr/.style={-{Stealth[scale=0.8]}},
  every node/.append style={align=center}
}
\usepackage{tcolorbox}
\tcbuselibrary{skins,breakable}
\usepackage{fontawesome}
\usepackage{fancyhdr}
\usepackage{enumitem}

\pagestyle{fancy}
\fancyhf{}
\renewcommand{\headrulewidth}{0.5pt}
\renewcommand{\footrulewidth}{0.5pt}
\fancyhead[L]{\small\bfseries\color{violet!85!black}\faBook\ SỔ TAY KIẾN THỨC TOÁN 2}
\fancyhead[R]{\small\bfseries\color{violet!85!black}Lớp Toán Thầy Lĩnh -- loptoanthaylinh.vn}
\fancyfoot[C]{\small\color{gray}-- Sổ tay kiến thức: Bài 20. Phép cộng có nhớ số có 2 chữ số với số có 2 chữ số --}

\setlength{\parindent}{0pt}
\setlength{\parskip}{0pt}

\begin{document}

% --- BANNER BÀI HỌC ---
\begin{tcolorbox}[
  enhanced,
  colback=violet!8!white,
  colframe=violet!70!black,
  arc=3mm, boxrule=1.2pt,
  left=3mm, right=3mm, top=2mm, bottom=2mm,
  halign=center
]
  {\LARGE\bfseries\color{violet!90!black}\faBookmark\ BÀI 20: PHÉP CỘNG (CÓ NHỚ) SỐ CÓ HAI CHỮ SỐ VỚI SỐ CÓ HAI CHỮ SỐ}\\[1mm]
  {\itshape\footnotesize\color{black!75} Kĩ thuật cộng hai số có 2 chữ số có nhớ sang hàng chục dạng 47 + 25 trên 1 trang A4}
\end{tcolorbox}

\vspace{-2mm}

\noindent
\begin{minipage}[t]{0.49\textwidth}\vspace{0pt}
\begin{tcolorbox}[
  enhanced,
  colback=violet!4!white,
  colframe=violet!75!black,
  arc=2mm, boxrule=0.8pt,
  left=2mm, right=2mm, top=1.8mm, bottom=1.8mm,
  title={\footnotesize\bfseries\color{white}\faKey\ 1. Quy tắc cộng 2 chữ số có nhớ}
]
{\footnotesize
\begin{itemize}\setlength{\itemsep}{1pt}\setlength{\leftskip}{-2.5mm}
  \item \textbf{Bước 1: Đặt tính thẳng cột:}\\
  Viết số hạng thứ hai dưới số hạng thứ nhất: hàng đơn vị thẳng hàng đơn vị, hàng chục thẳng hàng chục.
  \item \textbf{Bước 2: Cộng từ phải sang trái:}
  \begin{itemize}\setlength{\itemsep}{0.5pt}\setlength{\leftskip}{-2.5mm}
    \item \textbf{Hàng đơn vị:} Cộng hai chữ số đơn vị, viết chữ số đơn vị của tổng, \textbf{nhớ 1 sang hàng chục}.
    \item \textbf{Hàng chục:} Cộng hai chữ số hàng chục, rồi \textbf{cộng thêm 1 nhớ}, viết kết quả vào hàng chục.
  \end{itemize}
\end{itemize}
}
\end{tcolorbox}
\end{minipage}
\hfill
\begin{minipage}[t]{0.49\textwidth}\vspace{0pt}
\begin{tcolorbox}[
  enhanced,
  colback=teal!4!white,
  colframe=teal!75!black,
  arc=2mm, boxrule=0.8pt,
  left=2mm, right=2mm, top=1.8mm, bottom=1.8mm,
  title={\footnotesize\bfseries\color{white}\faSliders\ 2. Phân tích chi tiết ví dụ $47 + 25$}
]
{\footnotesize
\begin{center}
\begin{tabular}{r@{\,}ll}
  $4$ & $7$ & \quad \textbf{Thao tác tính:} \\
$+$ & & \quad $\bullet$ $7 + 5 = 12$, viết $\mathbf{2}$, \textbf{nhớ 1}. \\
  $2$ & $5$ & \quad $\bullet$ $4 + 2 = 6$, \\ \cline{1-2}
  $\mathbf{7}$ & $\mathbf{2}$ & \quad \ \ \textbf{thêm 1 bằng 7}, viết $\mathbf{7}$.
\end{tabular}
\end{center}
\vspace{1mm}
$\implies$ Kết quả: $47 + 25 = \mathbf{72}$.\\
\textbf{Công thức nhẩm hàng chục:}\\
$$\text{Chục kết quả} = a + c + 1$$
}
\end{tcolorbox}
\end{minipage}

\vspace{-1.5mm}
\begin{tcolorbox}[
  enhanced,
  colback=orange!4!white,
  colframe=orange!80!black,
  arc=2mm, boxrule=0.8pt,
  left=2mm, right=2mm, top=1.5mm, bottom=1.5mm,
  title={\footnotesize\bfseries\color{white}\faSitemap\ 3. Sơ đồ luồng xử lý phép cộng 2 chữ số có nhớ}
]
\centering
\begin{tikzpicture}[x=1.3cm, y=0.8cm]
  \node[rectangle, rounded corners=2pt, draw=violet!80!black, fill=violet!15, font=\scriptsize\bfseries] (dvi) at (0,0) {CỘNG HÀNG ĐƠN VỊ\\$7 + 5 = 12$\\Viết $\mathbf{2}$, nhớ $1$};
  \draw[->, very thick, red!85!black] (dvi) -- node[above, font=\tiny\bfseries] {Chuyển nhớ 1} (2.8,0);
  \node[rectangle, rounded corners=2pt, draw=teal!80!black, fill=teal!15, font=\scriptsize\bfseries] (chuc) at (4.5,0) {CỘNG HÀNG CHỤC\\$4 + 2 = 6$\\Thêm $1 = \mathbf{7}$};
  \draw[->, very thick, green!70!black] (chuc) -- node[above, font=\tiny\bfseries] {Ghép lại} (6.6,0);
  \node[rectangle, rounded corners=2pt, draw=green!80!black, fill=green!20, font=\normalsize\bfseries] at (7.4,0) {$\mathbf{72}$};
\end{tikzpicture}
\end{tcolorbox}

\vspace{-1.5mm}
\noindent
\begin{minipage}[t]{0.49\textwidth}\vspace{0pt}
\begin{tcolorbox}[
  enhanced,
  colback=green!4!white,
  colframe=green!60!black,
  arc=2mm, boxrule=0.8pt,
  left=2mm, right=2mm, top=1.8mm, bottom=1.8mm,
  title={\footnotesize\bfseries\color{white}\faCheckCircle\ 4. Ví dụ thực hành chuẩn mực}
]
{\footnotesize
\textbf{Đặt tính rồi tính:}\\
\begin{tabular}{r@{\,}lr@{\,}lr@{\,}l}
  $3$ & $6$ & $5$ & $8$ & $2$ & $9$ \\
$+$ & & $+$ & & $+$ & \\
  $2$ & $7$ & $1$ & $4$ & $4$ & $6$ \\ \hline
  $\mathbf{6}$ & $\mathbf{3}$ & $\mathbf{7}$ & $\mathbf{2}$ & $\mathbf{7}$ & $\mathbf{5}$
\end{tabular}
\quad
\begin{tabular}{l}
  $\bullet$ $6+7=13$, viết $3$ nhớ $1$; $3+2+1=6$.\\
  $\bullet$ $8+4=12$, viết $2$ nhớ $1$; $5+1+1=7$.\\
  $\bullet$ $9+6=15$, viết $5$ nhớ $1$; $2+4+1=7$.
\end{tabular}
}
\end{tcolorbox}
\end{minipage}
\hfill
\begin{minipage}[t]{0.49\textwidth}\vspace{0pt}
\begin{tcolorbox}[
  enhanced,
  colback=red!4!white,
  colframe=red!75!black,
  arc=2mm, boxrule=0.8pt,
  left=2mm, right=2mm, top=1.8mm, bottom=1.8mm,
  title={\footnotesize\bfseries\color{white}\faExclamationTriangle\ 5. Mẹo tránh bẫy trắc nghiệm}
]
{\footnotesize
\begin{itemize}\setlength{\itemsep}{1pt}\setlength{\leftskip}{-2.5mm}
  \item \textbf{Bẫy quên cộng 1 nhớ vào hàng chục:}\\
  $47 + 25 = 62$ (sai do chỉ lấy $4+2=6$ mà quên cộng thêm 1 nhớ). Kết quả đúng phải là $72$.
  \item \textbf{Mẹo ước lượng:}\\
  Vì có nhớ nên chữ số hàng chục của kết quả luôn lớn hơn tổng hai hàng chục đúng $1$ đơn vị.
\end{itemize}
}
\end{tcolorbox}
\end{minipage}


\end{document}
